\scanpage{114}

d'un él. de $\mathcal E_{X\deleted{\unsure{/Y}}}^T$, opérant
sur le sous-groupe $\mathcal E_{X/Y}^T$ par autom.
intérieurs.

Les notations \inserted{$\mathcal E_X$,} $\mathcal E_{X/Y}$ évoquent
intuitivement un $p^t$ de vue qui
regarde $\mathcal E_{X\deleted{\unsure{/Y}}}^T$ comme une \uline{extension}
de $\mathcal E_Y^T$ (ou d'un ss groupe
convenable seulement) par $\mathcal E_{X/Y}^T$,
qui sera souvent, pour les besoins de
la cause, noté $\unsure{\pi|\Pi}$ --- il joue le
rôle d'un groupe fondamental
(géométrique), tandis que $\mathcal E_Y^T$ joue
le rôle d'un groupe de Galois ---
ou d'un groupe profini mixte.

\margin{et j'hésite\newline
encore entre les\newline
notations $\mathcal E_{X/Y}^T$ et\newline
$\unsure{\pi|\Pi}_{X/Y}^T$; que je\newline
vais peut-être\newline
employer concur-\newline
remment.}

Le cas le plus important pour nous
sera celui où $Y=\operatorname{Spec}K$, $K$ un corps.
\deleted{\unsure{Aussi|Assez} \unsure{formellement|fréquemment|souvent}} Dans ce cas, $\mathcal E_Y\simeq\mathcal E_Y^{\unsure{\widetilde Y|\overline Y}}$

\margin{113}
